\documentclass[10pt,a4paper]{amsart}
\usepackage[margin=19mm]{geometry}
\usepackage{amsmath,amssymb,mathtools}
\usepackage[scr=rsfs,cal=euler]{mathalfa}
\usepackage{xfrac}
\usepackage[unicode,hidelinks,bookmarks=false]{hyperref}
\newcommand{\ie}{i.e.\ }
\newcommand{\cf}{cf.\ }
\newcommand{\resp}{respectively\ }
\newcommand{\Subsection}{\subsection}
\providecommand{\nobreakdash}{\nobreak\hbox{-}}
\setlength{\emergencystretch}{2em}
\multlinegap0pt
\usepackage{enumerate}
\usepackage{amssymb}
\newtheorem{lemma}{Lemma}
\title{Spectrally additive maps on the positive cones of the Wiener algebra}
\author{Shiho Oi, Kaito Sato}
\date{Source: arXiv:2512.01173v2 (CC BY 4.0)}
\begin{document}
\maketitle

\noindent\emph{Source and licence.} Spectrally additive maps on the positive cones of the Wiener algebra. Version arXiv:2512.01173v2 (CC BY 4.0), distributed by arXiv under the Creative Commons Attribution 4.0 International licence, http://creativecommons.org/licenses/by/4.0/, as recorded in the arXiv licence metadata for this exact version and shown on the abstract page https://arxiv.org/abs/2512.01173v2. Full licence notice: \texttt{licenses/arxiv-2512.01173-CC-BY-4.0.txt}.\par\smallskip

\noindent\textit{Editorial scope.} Counted unit: the article's lemma psi (source0.tex lines 153--155). The article's complete original proof of it is delivered with it (source0.tex lines 156--209, one \texttt{proof} environment of 54 lines). No mathematics is written, repaired or re-derived: the statement and the proof are the article's own lines, in the article's own notation, and no retained line is substituted or re-flowed.  Retained uncounted as the source-local context the proof needs: lines 82--84; lines 95--105; lines 138--140; lines 146--148.  Nothing was omitted from any retained span, so the package carries no omission record.  No retained line contains a citation, so the wrapper carries no bibliography and no bibliography entry.  No retained line contains a figure environment, an image-inclusion command or drawing code, so no image is delivered and the package declares no asset dependency.  Editorial formatting only, in addition: an AMS \texttt{amsart} preamble that restates the article's own declarations and macros used by the retained lines, namely the environments \texttt{lemma} on separate counters, together with the standard packages those lines need.  The retained environments print in this document as Lemma 1 in order of appearance, whereas the article prints the same items with the numbering of its own full document.  The complete native source of the article as distributed in the arXiv e-print for this exact version is bundled unchanged as \texttt{sources/190106/source0.tex}.
\begin{equation}\label{spectrumadditive}
\sigma(Tf+Tg)=\sigma(f+g) \quad  f, g \in A(\mathbb{T})_{+}.
\end{equation}

\begin{lemma}\label{2}
Let $m,n \in \mathbb{Z}$ with $m \neq n$. Then 
\[
M_{m,n}=\{ 2(e^{i\frac{2m\pi}{m-n}})^k \mid 0 \le k \le |m-n|-1\},
\]
and 
\[
\#(M_{m,n})=\frac{|m-n|}{\operatorname{gcd}(m,n)},
\]
where $\#(M_{m,n})$ is the cardinal number of the set $M_{m,n}$ and $\operatorname{gcd}(m,n)$ is the greatest common divisor of $m$ and $n$.
\end{lemma}

\begin{lemma}\label{B1}
The map $\psi:\mathbb{Z} \to \mathbb{Z}$ is a bijection.
\end{lemma}

\begin{lemma}\label{p0}
We have $\psi(0)=0$. 
\end{lemma}

\begin{lemma}\label{psi} 
We have either $\psi(n)=n$ for all $n \in \mathbb{Z}$, or $\psi(n)=-n$ for all $n \in \mathbb{Z}$. 
\end{lemma}

\begin{proof}
Let $m \in \mathbb{Z}\setminus \{1\}$. By \eqref{spectrumadditive}, we have 
\begin{equation}\label{spe1}
\sigma(z+z^{m})=\sigma(z^{\psi(1)}+z^{\psi(m)}).
\end{equation}
Hence we get $M_{1,m}=M_{\psi(1),\psi(m)}$. Lemma \ref{2} shows that $\frac{|m-1|}{\operatorname{gcd}(m,1)}=\frac{|\psi(m)-\psi(1)|}{\operatorname{gcd}(\psi(m),\psi(1))}$. Since $\psi(m) \neq \psi(1)$ by Lemma \ref{B1} and $\operatorname{gcd}(m,1)=1$, there exists $g \in \mathbb{Z}$ such that 
\begin{equation}\label{gcd1}
\psi(m)-\psi(1)=g(m-1).
\end{equation}
Note that $e^{in_1\theta}+e^{in_2\theta}=e^{i\frac{n_1+n_2}{2}\theta}(e^{i\frac{n_1-n_2}{2}\theta}+e^{i\frac{-n_1+n_2}{2}\theta})=2\cos(\frac{n_1-n_2}{2}\theta) e^{i\frac{n_1+n_2}{2}\theta}$. Let $r \in \mathbb{R}$ be an irrational number.  \eqref{spe1} implies that there is $\theta_0 \in [0, 2\pi]$ such that 
\[
2 \cos \left(\frac{m-1}{2} r\pi \right) e^{i\frac{m+1}{2}r \pi}=2 \cos \left(\frac{\psi(m)-\psi(1)}{2} \theta_0 \right) e^{i\frac{\psi(m)+\psi(1)}{2}\theta_0}.
\]
Thus we have
\begin{equation}\label{absolutevalue}
\left|\cos \frac{m-1}{2} r\pi \right|=\left| \cos \frac{\psi(m)-\psi(1)}{2} \theta_0 \right|
\end{equation}
and 
\begin{equation}\label{argument}
e^{i(\frac{\psi(m)+\psi(1)}{2}\theta_0-\frac{m+1}{2}r \pi)} \in \mathbb{R}.
\end{equation}
According to \eqref{absolutevalue},  we have two following cases:\\
{\it{First case} }: $\cos \frac{m-1}{2} r\pi = \cos \frac{\psi(m)-\psi(1)}{2} \theta_0$. Then for some $n \in \mathbb{Z}$,
$\frac{m-1}{2} r\pi =\pm \frac{\psi(m)-\psi(1)}{2} \theta_0+2n\pi$.
By \eqref{gcd1}, we have $\frac{m-1}{2} r\pi =\pm \frac{g(m-1)}{2} \theta_0+2n\pi$, so that $(m-1)(r \pi \mp g\theta_0)=4n\pi$. Since $\mathbb{Z} \ni m \neq 1$, it follows that there is $q \in \mathbb{Q}$ such that 
$g \theta_0=r\pi + q\pi$ or $g \theta_0=-r\pi + q\pi$. \\
{\it{Second case:}} $-\cos \frac{m-1}{2} r\pi = \cos \frac{\psi(m)-\psi(1)}{2} \theta_0$, which is equivalent to $\cos ( \pi - \frac{m-1}{2} r\pi)= \cos \frac{\psi(m)-\psi(1)}{2} \theta_0$. Then there is $n \in \mathbb{Z}$ such that 
$\frac{\psi(m)-\psi(1)}{2} \theta_0=\pm (\pi - \frac{m-1}{2} r\pi)+2n \pi$.
By a similar argument as in the first case, we also obtain $q \in \mathbb{Q}$ such that 
$g \theta_0=r\pi + q\pi$ or $g \theta_0=-r\pi + q\pi$. \\ 
Ultimately, we conclude that \eqref{absolutevalue} implies that there exists $q \in \mathbb{Q}$ such that 
\begin{equation}\label{gt}
g \theta_0=r\pi + q\pi \quad \text{or}  \quad g \theta_0=-r\pi + q\pi.
\end{equation}
 By \eqref{argument}, there is $n \in \mathbb{Z}$ such that 
\[
\frac{\psi(m)+\psi(1)}{2}\theta_0-\frac{m+1}{2}r \pi=n \pi.
\]
Multiplying both sides by $g$ and substituting \eqref{gt}, we obtain 
\[
((m+1)g \mp (\psi(m)+\psi(1)))r =(\psi(m)+\psi(1))q -2ng.
\]
Note that $(m+1)g \mp (\psi(m)+\psi(1))$ is an integer, $r$ is irrational, and the right-hand side is a rational number. Thus we obtain $\psi(m)+\psi(1)=\pm (m+1)g$. When $\psi(m)+\psi(1)=(m+1)g$ (resp.  $\psi(m)+\psi(1)=-(m+1)g$), it follows from \eqref{gcd1} that
$\psi(m)=\psi(1)m$ (resp. $\psi(1)=\psi(m)m$).  Thus it follows that for any $m \in \mathbb{Z} \setminus \{1\} $, 
\[
\psi(m)=\psi(1)m  \quad \text{or} \quad \psi(1)=\psi(m)m
\]
holds. Let $m \in \mathbb{Z}$ with $|m|>|\psi(1)|$. Suppose $\psi(1)=\psi(m)m$, then $\psi(m)=\frac{\psi(1)}{m}$. This contradicts that  $\psi(m) \in \mathbb{Z}$. Thus if  $|m|>|\psi(1)|$, $\psi(m)=\psi(1)m$ holds.  Combining this with Lemma \ref{B1},  we obtain $\psi(1)=1$ or $\psi(1)=-1$. 
Therefore if  $|m|>1$, then 
\begin{equation}\label{eq}
\psi(m)=\psi(1)m.
\end{equation}
By Lemma \ref{p0}, \eqref{eq} holds for $m=0,1$.  Finally, since $\psi$ is a bijection, \eqref{eq}  also holds for $m=-1$. Hence we obtain the desired result.
\end{proof}

\end{document}
